\section{Ranked candidate gaps}
\label{sec:gaps}

\subsection{How a candidate becomes a thesis gap}

An unanswered question is not automatically a useful contribution.  In
this review, a candidate is retained only when it satisfies four
conditions:

\begin{enumerate}
  \item it is not answered by the supplied corpus;
  \item CRESCO8 provides an unusual or well-controlled setting in which
        to answer it;
  \item the causal variables and outcome can be measured with available
        tools; and
  \item a result in either direction would still be informative.
\end{enumerate}

The ranking below is deliberately conservative.  ``Not answered by the
corpus'' is the current evidence statement.  ``Novel in the field''
remains to be established by a broader literature search before the
thesis proposal is frozen.

\subsection{Priority 1: controlled memory-tier evidence}
\label{sec:gap-primary}

\paragraph{\researchq{1}.}
\emph{For quantized CPU-only LLM inference on one Xeon Max~9480 socket,
how do verified local-HBM and local-DDR placement affect prefill and
sustained decode under otherwise identical conditions?}

This is the primary gap.  Na et al.\ measure LLM inference on Xeon Max
but do not isolate memory tier or quantized \lcpp{} execution.  Ibeid et
al.\ isolate memory behaviour but not LLM inference.  Goto et al.\ stage
selected HPL data in HBM but do not report an isolated HBM effect.
The experiment therefore connects evidence that is currently separate.

The comparison has two pre-registered hypotheses:

\begin{itemize}
  \item \(H_{1a}\): when the measured live set fits in local HBM,
        sustained decode obtains a larger relative improvement than
        prefill because its low-batch matrix--vector work has lower
        arithmetic intensity.
  \item \(H_{1b}\): the observed performance ratio is smaller than the
        local STREAM bandwidth ratio because not all inference time is
        spent streaming the relevant data.
\end{itemize}

Both hypotheses may be rejected without invalidating the study.  A
negligible HBM effect accompanied by low memory bandwidth would show
that the chosen phase is not bandwidth-bound.  A negligible effect
despite high bandwidth pressure would direct attention to cache reuse,
kernel coverage, synchronization or placement errors.  A result close
to the STREAM ratio would be surprising and would require especially
careful validation.

The minimum sufficient contribution is a phase-resolved comparison with
raw data, uncertainty, live-set accounting, page residency and achieved
bandwidth.  It requires no runtime patch and should be completed before
any secondary gap.

\subsection{Priority 2: the capacity and KV boundary}
\label{sec:gap-capacity}

\paragraph{\researchq{2}.}
\emph{How does the HBM benefit change as the verified live set
approaches and crosses usable local-HBM capacity, and which component
of the live set causes the transition?}

This question distinguishes at least three regimes:

\begin{enumerate}
  \item weights, KV state and transient allocations fit in local HBM;
  \item weights fit but context, batch or packed buffers push the full
        live set beyond the tier; and
  \item weights alone exceed usable local HBM.
\end{enumerate}

The relevant independent variables are model/quantization, prompt
length, generated length, batch or concurrent sequences, and KV
format.  The dependent variables remain TTFT, ITL, phase throughput and
achieved bandwidth, with peak resident memory and tier-resolved pages
as mechanism checks.

The main threat is mistaking an allocation failure, paging event or
direct-mapped Cache-mode conflict for a smooth capacity effect.
Flat-mode runs should therefore establish the primary boundary.
Cache mode is a separate policy experiment, not a replacement for
verified Flat placement.

This question can be answered in a useful first form with whole-process
placement: compare all-HBM runs that fit against matched DDR runs and
document where the HBM run ceases to be feasible.  Independent weight
and KV placement is needed only for the stronger question of how to
operate beyond that boundary.

\subsection{Priority 3: topology as a controlled factor}
\label{sec:gap-topology}

\paragraph{\researchq{3}.}
\emph{When additional sockets or SNC domains are used, can explicit
execution and page placement avoid the regressions observed with an
unmanaged configuration?}

The corpus motivates this question through a tension: Na et al.\ report
poor naive 96-core and SNC4 inference results, while Ibeid et al.\ and
Goto et al.\ benefit from explicit domain-local mapping.  A useful
experiment should compare named policies rather than repeat a generic
``one versus two socket'' sweep:

\begin{itemize}
  \item one socket and its local memory;
  \item two independent, socket-local replicas, reporting aggregate and
        per-request performance;
  \item one process spanning sockets without deliberate placement, as a
        negative control; and
  \item a partitioned design only if the runtime can express it without
        changing unrelated kernels.
\end{itemize}

SNC4 should be attempted only if the administrators can expose and
support the mode for allocated nodes.  Remote page counts, UPI traffic
and per-domain bandwidth are required to interpret the result.  If
mode changes are unavailable, the thesis should report that operational
constraint and retain socket-local placement as the topology study.

This gap is secondary because a two-socket result is difficult to
explain before the one-socket tier effect is established.

\subsection{Conditional software gap: selective tiering}
\label{sec:gap-software}

\paragraph{\researchq{4}.}
\emph{When the live set exceeds HBM, can an explicit CPU-side placement
of weights and KV state improve performance over the best unmodified
whole-process policy?}

The corrected Aurora paper makes this a plausible systems question:
its bulk matrix resides in DDR and selected panel data are staged into
HBM~\cite{goto2026aurorahpl}.  Transformer data have different access
patterns, so the useful split must be measured rather than copied.  Two
candidate policies are:

\begin{itemize}
  \item weights in HBM and KV/scratch in DDR, prioritising the repeated
        weight stream in low-batch decode; and
  \item KV state in HBM and weights in DDR when long context makes
        attention-state traffic dominant.
\end{itemize}

The static audit shows that this policy is not exposed as an ordinary
CPU-side \lcpp{} control at the pinned commit.  That creates a possible
software contribution, but only after \researchq{1} and \researchq{2} demonstrate a
benefit and a capacity constraint.  The contribution would consist of
the smallest auditable control that expresses the policy, not a general
NUMA scheduler.

The implementation succeeds only if it beats the strongest existing
baseline, not merely an unmanaged run.  Required ablations are:
all-DDR, all-HBM where feasible, operating-system interleave or preferred
policy if supported, and each selective placement.  Page residency and
allocation failures must be reported.

\subsection{Extensions, not core gaps}

\paragraph{Energy.}
Energy per token is useful only if the accessible interface has a clear
scope, resolution and sampling method.  A node-level measurement cannot
be presented as CPU energy without accounting for the rest of the node.
Energy should therefore follow, rather than delay, the validated
performance study.

\paragraph{Online serving and tail latency.}
Arrival processes, scheduling and request interference matter for
deployment, but they introduce a second class of software effects.
They are appropriate after the single-request mechanism and capacity
boundary are understood.

\paragraph{Cache mode.}
Cache mode may provide a no-code overflow policy, but direct-mapped
conflicts and undocumented firmware policy can complicate causal
interpretation.  It is a valuable extension if mode changes are
available, not a prerequisite.

\paragraph{GPU and multi-node comparisons.}
Neither is needed for the memory-tier claim.  A fair GPU serving
comparison would require matched quality, software maturity, batching,
capacity and power boundaries.  Multi-node LLM inference adds network
and parallelisation variables not addressed by the primary question.

\subsection{Priority and dependencies}

This ordering turns the original catalogue of many speculative gaps
into one primary question, two evidence-driven follow-ons and one
conditional implementation.  It also preserves a complete thesis if
the runtime never needs to be changed.

\begin{table}[htbp]
\centering
\small
\caption{Recommended priority.  A later item starts only when the
preceding evidence justifies it.}
\label{tab:gap-priority}
\begin{tabularx}{\textwidth}{@{}P{1.25cm}P{3.35cm}P{2.65cm}YY@{}}
\toprule
\textbf{Order} & \textbf{Question} & \textbf{Role} &
\textbf{Entry condition} & \textbf{Useful result if null} \\
\midrule
1 & \researchq{1}: local HBM vs local DDR & Core empirical contribution
& Platform and placement validation complete
& Quantifies where HBM does not matter and identifies the limiting
  resource. \\
\addlinespace
2 & \researchq{2}: capacity/KV boundary & Core or secondary empirical contribution
& At least one stable configuration and a measured live-set model
& Establishes the feasible envelope even if no sharp transition is
  observed. \\
\addlinespace
3 & \researchq{3}: topology-aware scaling & Secondary
& Remote-traffic counters and supported topology controls
& Shows whether extra resources are neutral or harmful under explicit
  policies. \\
\addlinespace
4 & \researchq{4}: selective tiering & Conditional systems contribution
& Reproducible HBM benefit, capacity pressure and missing existing control
& A negative result bounds the value of fine-grained placement. \\
\addlinespace
5 & Energy, serving, Cache mode & Extensions
& Validated measurement interfaces and remaining time
& Adds deployment context without being required for thesis completeness. \\
\bottomrule
\end{tabularx}
\end{table}
